(1) 写$\mathcal A E=\omega^2\epsilon E$，其中$\mathcal A=\nabla\times\mu^{-1}\nabla\times$在理想导体边界下为自伴算符。对$\epsilon\to\epsilon+\delta\epsilon$和$\omega\to\omega+\delta\omega$作一阶变分，左乘未扰动场$E_0^*$积分，利用自伴性消去$\delta E$项，得$2\omega\delta\omega\int\epsilon|E_0|^2dV=-\omega^2\int\delta\epsilon|E_0|^2dV$。所以$\delta\omega/\omega=-\frac12[\int\delta\epsilon|E_0|^2dV]/[\int\epsilon|E_0|^2dV]$。
(2) 设原介质$\epsilon$均匀，横截面积$A$，薄片厚$\ell\ll L$且位于$x=x_0$，其介电增量为常数$\Delta\epsilon$。分母为$\epsilon A E_a^2\int_0^L\sin^2(\pi x/L)dx=\epsilon A E_a^2L/2$；分子主阶约$\Delta\epsilon A E_a^2\ell\sin^2(\pi x_0/L)$，故$\delta\omega/\omega\simeq-(\Delta\epsilon/\epsilon)(\ell/L)\sin^2(\pi x_0/L)$。
(3) 对$\Delta\epsilon>0$，频移恒为非正：插入高介电材料会降低谐振频率。场节点$x_0=0,L$附近主阶频移近零，反节点$x_0=L/2$最大，值为$-(\Delta\epsilon/\epsilon)(\ell/L)$。该位置依赖可用于腔场测绘或介质传感，且对归一化振幅$E_a$不敏感。
(4) 若模式简并，必须在简并模子空间内对$\delta\epsilon$的矩阵元作对角化；若介电常数随频率变化，正确归一化涉及能量密度的$\partial_\omega(\omega\epsilon)$，不能直接套用非色散分母。若$\Delta\epsilon$或片厚过大，场本身会明显重排，需要求解完整边值问题；小扰动公式只提供首阶渐近。
